\documentclass[11pt,a4paper]{article}
\usepackage[margin=27mm]{geometry}
\usepackage{amsmath,amssymb,amsthm,fontspec,unicode-math,hyperref,booktabs,longtable}
\setmainfont{STIXGeneral}
\setmathfont{STIX Two Math}
\setmonofont{Menlo}
\usepackage{newunicodechar}
\newunicodechar{∎}{\ensuremath{\blacksquare}}
\hypersetup{hidelinks}
\setlength{\emergencystretch}{2em}
\providecommand{\tightlist}{\setlength{\itemsep}{0pt}\setlength{\parskip}{0pt}}
\newcommand{\gen}[1]{\langle #1\rangle}
\newcommand{\Sym}{\operatorname{Sym}}
\newcommand{\Alt}{\operatorname{Alt}}
\newcommand{\core}{\operatorname{core}}
\newcommand{\F}{\mathbb F}
\newtheorem{theorem}{Theorem}
\newtheorem{lemma}[theorem]{Lemma}
\newtheorem{proposition}[theorem]{Proposition}
\newtheorem{corollary}[theorem]{Corollary}

\title{Disjoint-subset geometries and their full graph automorphisms}
\author{Human Compression project}
\date{}
\begin{document}
\maketitle
\begin{abstract}
A weighted star constructs the complete family of disjoint-subset geometries. We determine its residues, panel sizes, distances, gonality and diagram, and reconstruct both the ground-set action and the type classes from the uncolored incidence graph. The spectral argument is included in full, so the full automorphism claim does not follow merely from its type-preserving counterpart. The thin boundary, the diagram order and the locally Desargues group are stated with their reviewed corrections.
\end{abstract}
\section{Introduction}
Garcia-Colin and Leemans constructed locally Kneser graphs through incidence geometries of arbitrary rank. Their construction is a literal specialization of the forest geometry [UB], not only a similarity of diagrams. That shared theorem supplies flag completion, residual connectedness and the full residue structure. The distance formulas and both reconstruction arguments are local to this application and are proved below. We use the ten-page institutional revised author edition, with its own title and theorem numbering; no identity with a publisher or earlier arXiv edition is presumed.
\section{The weighted-star construction}
Let \(r\geq2\), \(k\geq1\), and \(N>rk\) be integers. On a star with \(r\) labelled leaves give each leaf a block of size \(k\), and give the central vertex a block of size \(N-rk\). Root the star at its center. Its edge cuts are precisely the leaf blocks. Theorem UB.F therefore gives a geometry \(\mathcal K(N,k,r)\) whose type set is \(\{1,\ldots,r\}\), with \begin{equation}
\{\text{type-}i\text{ elements}\}=\binom{[N]}k,
\qquad (i,A)*(j,B)\iff i\ne j\text{ and }A\cap B=\varnothing
\label{eq:disjoint-subset-geometries:7-1}
\end{equation} for distinct elements. This is exactly Construction 1.1 of {[}GL21{]}, on setting \begin{equation}
N=n+k(r-2),\qquad n\geq2k+1.
\label{eq:disjoint-subset-geometries:7-2}
\end{equation} The case \(k=1,n=3\) is the unit-star geometry of Theorem UB.F, so the two constructions overlap at the level of their actual cosets and actions.

\textbf{Theorem B4.Star (the weighted-star specialization).} The geometry \(\mathcal K(N,k,r)\) is firm, residually connected and flag-transitive under \(S_N\). A flag of \(s\) elements has residue \(\mathcal K(N-sk,k,r-s)\), on the complement of its subsets, with the evident rank-zero and rank-one conventions. Each type has \(\binom Nk\) elements, and each panel has \begin{equation}
\binom{N-(r-1)k}{k}=\binom{n-k}{k}
\label{eq:disjoint-subset-geometries:7-3}
\end{equation} elements. It is thick except when \(k=1,n=3\), when it is thin. The incidence graph is \(KG(N,k)\otimes K_r\); its vertex links are \(KG(N-k,k)\otimes K_{r-1}\). The group \(S_N\) is transitive on its maximal cliques.

\emph{Proof.} All positive-weight and laminar-cut hypotheses of UB.F hold for the displayed star, so firmness, residual connectedness, flag completion and flag transitivity follow from that theorem. Explicitly, a flag is a family of disjoint \(k\)-subsets assigned to distinct types. Its residue uses the unused points and types, proving the residue identity. It also gives \eqref{eq:disjoint-subset-geometries:7-3}. If \(k=1\), this number is \(n-1\), which is two only for \(n=3\) and is at least three otherwise. If \(k\geq2\), then \(n-k\geq k+1\), so it is at least \(\binom{k+1}k=k+1\geq3\). This proves the exact thin boundary without an additional parameter restriction.

Equation \eqref{eq:disjoint-subset-geometries:7-1} is the tensor-product adjacency relation. Fixing \((i,A)\) leaves the other \(r-1\) types and the \(N-k\) points outside \(A\), proving the link formula. Every maximal clique is a chamber by flag completion. A bijection between the \(r\) disjoint subsets of two chambers extends to a permutation of their complements, proving the stated transitivity. ∎

Here \(\otimes\) denotes the categorical, or tensor, graph product: both coordinates must be adjacent. For \(r=2\), its link formula is simply an edgeless graph on the neighbors, as \(K_1\) has no edges. The rank-two residues are the bipartite doubles of \(KG(n,k)\). We compute their distances without invoking a classification of distance-regular graphs.

\textbf{Proposition B4.Distance (replacement distances).} Put \(s=n-2k>0\), and let \(A,B\) be \(k\)-subsets of \([n]\), with \(c=|A\cap B|\). In the bipartite double of \(KG(n,k)\), the distances between copies in the same and in opposite parts are respectively \begin{equation}
d_{\mathrm{even}}(A,B)=2\left\lceil\frac{k-c}{s}\right\rceil,
\qquad
d_{\mathrm{odd}}(A,B)=2\left\lceil\frac{c}{s}\right\rceil+1.
\label{eq:disjoint-subset-geometries:7-4}
\end{equation} Its diameter, from either part, is \(2\lceil k/s\rceil+1\). Its gonality is three for \(n=2k+1\) and two for \(n\geq2k+2\). The odd girth of \(KG(n,k)\) is \(2\lceil k/s\rceil+1\).

\emph{Proof.} If \(A,C,D\) are successive vertices of a disjointness walk, then \(A\cup D\) lies in the complement of \(C\), of size \(n-k=k+s\). Thus \(|A\setminus D|\leq s\): two steps change at most \(s\) entries. Conversely, if two \(k\)-sets differ in at most \(s\) entries, their union has size at most \(k+s\), leaving at least \(k\) points from which to choose a common disjoint set. Every such replacement is realizable in two steps.

An even walk from \(A\) to \(B\) must replace the \(k-c\) elements of \(A\setminus B\), so needs at least \(2\lceil(k-c)/s\rceil\) steps. Partition those elements and the corresponding new elements of \(B\setminus A\) into matching batches of size at most \(s\). Replace one batch at a time and use the common disjoint set just constructed. This realizes the bound, including distance zero when \(A=B\).

For an odd walk, let \(C\) be its penultimate set; then \(C\cap B=\varnothing\), so \(|A\cap C|\leq k-c\). Its even initial segment has length at least \(2\lceil c/s\rceil\). Equality is achievable: take all of \(A\setminus B\), and fill \(C\) with \(c\) points outside \(A\cup B\). There are \(s+c\) available such points. Use the even construction from \(A\) to \(C\), then the edge to \(B\). This proves \eqref{eq:disjoint-subset-geometries:7-4}. Both parity distances are finite, so the double is connected. Their largest value is the asserted odd distance between the two copies of a single set.

The number of common neighbors of two distinct sets with intersection size \(c\) is \begin{equation}
\binom{n-2k+c}{k},
\label{eq:disjoint-subset-geometries:7-5}
\end{equation} with the binomial coefficient zero if its upper argument is smaller than \(k\). For \(n=2k+1\), distinct sets have \(c\leq k-1\), so \eqref{eq:disjoint-subset-geometries:7-5} is at most one. There is no four-cycle. To exhibit a six-cycle, take three \(k\)-sets \(S\cup\{x\},S\cup\{y\},S\cup\{z\}\), where \(|S|=k-1\) and \(x,y,z\) are distinct outside \(S\). In the other part take the complements of the unions of each consecutive pair. Each complement has size \(k\); the three complements are distinct, and they yield a six-cycle. This works also for \(k=1\). If \(n\geq2k+2\), choose two sets differing in one entry. Formula \eqref{eq:disjoint-subset-geometries:7-5} is at least \(\binom{k+1}k\geq2\), giving a four-cycle. These facts give the gonalities.

Finally, \eqref{eq:disjoint-subset-geometries:7-4} with \(A=B\) gives the minimum length of an odd closed walk in the Kneser graph. Every odd closed walk contains an odd cycle: if it repeats a vertex internally, splitting it at that vertex leaves an odd closed subwalk, and induction removes repetitions. Therefore the minimum odd closed walk is itself an odd cycle, proving the odd-girth assertion. ∎

It follows that the Buekenhout diagram is complete, with each edge labelled \(d-g-d\), where \(d=2\lceil k/(n-2k)\rceil+1\) and \(g\) is as in B4.Distance. No rank-two residue is a generalized digon: the two copies of the same \(k\)-set are not incident. Thus the geometric property \((IP)_2\) holds exactly for \(n=2k+1\). The order at each diagram vertex is the panel size minus one, namely \(\binom{n-k}{k}-1\).
\section{Recovering the permutation action}
\textbf{Lemma B4.Johnson.} If \(N>2k\), every automorphism of the Johnson graph \(J(N,k)\) is induced by a unique permutation of \([N]\).

\emph{Proof.} The vertices are the \(k\)-sets, with adjacency when their intersection has size \(k-1\). For \(k=1\) this is the complete graph on \([N]\), so the assertion holds. Assume \(k\geq2\). For each \((k-1)\)-set \(C\), the family of \(k\)-sets containing \(C\) is a clique, called its star. For each \((k+1)\)-set \(D\), the family of \(k\)-sets contained in \(D\) is a clique, called its top.

These are all maximal cliques. To see this, take two adjacent vertices \(A=C\cup\{a\}\), \(B=C\cup\{b\}\). Any further vertex adjacent to both either contains \(C\), or has the form \((C\setminus\{x\})\cup\{a,b\}\) for some \(x\in C\). In the latter case it is contained in \(D=C\cup\{a,b\}\). A vertex \(C\cup\{z\}\) with \(z\notin D\) meets such a latter vertex in only \(k-2\) points, so cannot occur in the same clique. Consequently a clique through \(A,B\) lies entirely in their star or their top, and maximality gives the whole star or top. Both displayed families are maximal by the same observation.

Stars have size \(N-k+1\), larger than the top size \(k+1\); hence every automorphism preserves stars. Two distinct stars intersect if and only if their \((k-1)\)-sets differ in one point, so the induced action on stars is an automorphism of \(J(N,k-1)\). Induction identifies it with a permutation of \([N]\), since \(N>2(k-1)\). A \(k\)-set is determined by the stars containing it, namely its \((k-1)\)-subsets. The action on stars therefore determines the original action, proving existence and uniqueness. ∎

\textbf{Proposition B4.TypedAut.} The type-preserving automorphism group of \(\mathcal K(N,k,r)\) is \(S_N\), and the group allowed to permute its types is \(S_N\times S_r\).

\emph{Proof.} Within one type, count the common neighbors of two distinct vertices in a fixed other type. Formula \eqref{eq:disjoint-subset-geometries:7-5}, now with ambient size \(N\), is strictly largest exactly when their intersection has size \(k-1\): its upper argument there is \(N-k-1\geq k\), and decreasing the intersection either strictly decreases the binomial coefficient or makes it zero. Thus the incidence relation reconstructs the Johnson adjacency in each type. By B4.Johnson a type-preserving automorphism acts on type \(i\) by a permutation \(\pi_i\) of \([N]\).

These permutations are all equal. For two types \(i\ne j\), compose with \(\pi_i^{-1}\) on the ground set. The resulting permutation fixes every vertex of type \(i\), and sends a type-\(j\) set \(B\) to \(\pi_i^{-1}\pi_j B\) with the same disjointness neighborhood. Distinct \(k\)-sets have distinct such neighborhoods: if \(x\in B\setminus B'\), choose a \(k\)-set containing \(x\) in the complement of \(B'\); it meets \(B\) but is disjoint from \(B'\). This is possible because \(N-k\geq k\). Hence every \(B\) is fixed, and faithfulness of the action on \(k\)-sets gives \(\pi_i=\pi_j\). The diagonal \(S_N\) acts, so it is exactly the type-preserving group.

Every permutation of the \(r\) types acts independently and commutes with the diagonal \(S_N\). Composing any type-permuting automorphism with the inverse of its type permutation reduces to the type-preserving case. The two actions have trivial intersection, giving the direct product. ∎
\section{Recovering the types from the uncolored graph}
The automorphism group of a typed geometry need not be the full group of its incidence graph: a graph automorphism could conceivably mix the type classes. We exclude this possibility for \eqref{eq:disjoint-subset-geometries:7-1} by identifying the type classes intrinsically as the maximum independent sets.

Let \(\mathscr V=\binom{[N]}k\), \(v=|\mathscr V|\), and let \(K\) be its disjointness matrix. Its row sum is \(d=\binom{N-k}{k}\). For a subset \(T\) of size \(j\leq k\), put \(f_T(A)=1\) if \(T\subseteq A\), and zero otherwise. Let \(V_j\) be the span of the \(f_T\) with \(|T|=j\). The spaces form an increasing filtration: for \(|S|=j-1\), \begin{equation}
\sum_{\substack{T\supset S\\ |T|=j}} f_T=(k-j+1)f_S.
\label{eq:disjoint-subset-geometries:7-6}
\end{equation} Also \(V_0\) consists of constants and \(V_k=\mathbb R^{\mathscr V}\), because the latter indicators are the individual vertex indicators.

For \(|T|=j\), direct counting and inclusion--exclusion give \begin{equation}
(Kf_T)(A)=\binom{N-k-j}{k-j}\,1_{A\cap T=\varnothing}
=\binom{N-k-j}{k-j}\sum_{S\subseteq T}(-1)^{|S|}f_S(A).
\label{eq:disjoint-subset-geometries:7-7}
\end{equation} Thus \(K\) preserves the filtration and acts on \(V_j/V_{j-1}\) by \begin{equation}
\lambda_j=(-1)^j\binom{N-k-j}{k-j}.
\label{eq:disjoint-subset-geometries:7-8}
\end{equation} Choosing bases adapted to the filtration shows that every eigenvalue is one of these values; any zero-dimensional quotient simply contributes none. As \(K\) is real symmetric, it is diagonalizable. For \(j\geq1\), \begin{equation}
|\lambda_j|\leq\binom{N-k-1}{k-1}
=\frac{k}{N-k}\,d<d.
\label{eq:disjoint-subset-geometries:7-9}
\end{equation} The inequality follows successively from \(\binom{a-1}{b-1}/\binom ab=b/a<1\). In particular \(d\) occurs only on the constant subspace.

The graph adjacency matrix is \(M=(J_r-I_r)\otimes K\). The two eigenvalues of \(J_r-I_r\) are \(r-1\) on constants and \(-1\) on their orthogonal complement. Tensor products of eigenbases give all eigenvectors of \(M\). Since \(N>rk\), \begin{equation}
(r-1)\frac{k}{N-k}<1.
\label{eq:disjoint-subset-geometries:7-10}
\end{equation} It follows from \eqref{eq:disjoint-subset-geometries:7-9} that the least eigenvalue of \(M\) is exactly \(-d\). Its eigenspace consists precisely of vectors constant on each type, with the \(r\) constants summing to zero.

\textbf{Theorem B4.Aut (the full incidence-graph group).} The maximum independent sets of the graph of \(\mathcal K(N,k,r)\) are exactly its \(r\) type classes. Its full graph automorphism group is \(S_N\times S_r\).

\emph{Proof.} The graph has \(rv\) vertices and valency \(t=(r-1)d\). If an independent set has size \(a\), write its indicator as \(x=\alpha\mathbf1+z\), where \(\alpha=a/(rv)\) and \(z\perp\mathbf1\). Then \(\|z\|^2=rv\alpha(1-\alpha)\). Independence gives \(x^TMx=0\), hence \[
z^TMz=-t\,rv\alpha^2.
\] The least-eigenvalue inequality \(z^TMz\geq-d\|z\|^2\) therefore gives \(t\alpha\leq d(1-\alpha)\), or \(\alpha\leq1/r\). Thus \(a\leq v\), attained by each type class.

If \(a=v\), equality holds in the eigenvalue inequality. Expanding \(z\) in an orthonormal eigenbasis shows that every component on a larger eigenvalue is zero. Thus \(z\), and therefore \(x\), is constant on each type. Since \(x\) is a zero-one indicator of size exactly \(v\), it is the indicator of one type class. Graph automorphisms preserve maximum independent sets and hence permute the types. Proposition B4.TypedAut now gives the full group. ∎

For \(k=2,n=5,r=3\), the graph has \(3\binom72=63\) vertices and each link is the bipartite double of the Petersen graph \(KG(5,2)\), the Desargues incidence graph. Explicitly, complementing the sets in one part identifies its incidence with containment between the two-subsets and three-subsets of a five-element set, the classical Desargues \(10_3\) configuration. Its full group is \(S_7\times S_3\), of order \(30240\). The source's displayed Corollary 5.1 names only \(S_7\); that is the type-preserving subgroup. Theorem B4.Aut establishes the larger full graph group without relying on a finite enumeration.
\section{The corrected source conclusions}
\textbf{Corollary B4.Source (disjoint-subset geometries and locally Kneser graphs).} For \(r\geq2\), \(k\geq1\), \(n\geq2k+1\), the geometry of {[}GL21, Construction 1.1{]} is the weighted star \eqref{eq:disjoint-subset-geometries:7-1}--\eqref{eq:disjoint-subset-geometries:7-2}. It has all the existence, residue and clique-transitivity properties in B4.Star, the complete diagram and distance parameters in B4.Distance, and the type-preserving and full graph groups in B4.TypedAut and B4.Aut. It satisfies \((IP)_2\) exactly when \(n=2k+1\).

\emph{Proof.} The star has positive central weight \(n-2k\), so UB.F gives B4.Star. Its rank-two residues are the bipartite doubles of B4.Distance; they are never digons, and their four-cycle criterion gives \((IP)_2\). B4.TypedAut and B4.Aut prove the two action assertions and exclude extra uncolored graph automorphisms. These establish {[}GL21, Theorem 1.2 and Corollaries 1.3, 1.4, 4.8, 5.1{]} with three explicit corrections: the diagram order is \(\binom{n-k}{k}-1\); the case \(k=1,n=3\) is thin while every other allowed case is thick; and the locally Desargues example has full group \(S_7\times S_3\). ∎

\section*{References and retained inputs}
{[}UB{]} \emph{Forest coordinates and relative generation}, common theory in the present volume. Its full proofs and definitions are part of this application's mathematical dependencies.

{[}GL21{]} N. Garcia-Colin and D. Leemans, \emph{An infinite family of locally X graphs based on incidence geometries}, Art of Discrete and Applied Mathematics \textbf{4} (2021), \#P3.01, DOI 10.26493/2590-9770.1328.902. The consulted institutional revised author manuscript is titled \emph{An infinite family of incidence geometries whose incidence graphs are locally X}. Its Theorem 1.2 and Corollaries 1.3, 1.4, 4.8 and 5.1 provide the numbering used here; the earlier arXiv version has different numbering. No identity of the institutional PDF with the publisher's typeset text is presumed.

\end{document}
